% part: intuitionistic-logic
% chapter: semantics
% section: topological-semantics

\documentclass[../../../include/open-logic-section]{subfiles}

\begin{document}

\olfileid{int}{sem}{top}

\olsection{सांस्थितिक चिन्हार्थमीमांसा}

अंतःप्रज्ञावादी तर्कशास्त्राला चिन्हार्थमीमांसा
देण्याचा आणखी एक मार्ग म्हणजे संस्थिति ही
गणिती संकल्पना वापरणे.

\begin{defn}
  $X$ हा संच असू द्या. $X$ वरील
  \emph{संस्थिति} म्हणजे खालील गुणधर्म पूर्ण
  करणारा संच $\Top{O}
  \subseteq \Pow{X}$. $\Top{O}$ च्या
  !!{element}s ना त्या संस्थितीचे
  \emph{विवृत संच} म्हणतात. $\Top{O}$
  सहित~$X$ या संचाला \emph{सांस्थितिक अवकाश}
  म्हणतात.
  \begin{enumerate}
  \item रिकामा संच आणि संपूर्ण अवकाश विवृत आहेत:
    $\emptyset$, $X \in \Top{O}$.
  \item विवृत संच सांत छेदांखाली संवृत असतात:
    $U$, $V \in \Top{O}$ असेल, तर
    $U \cap V \in \Top{O}$.
  \item विवृत संच स्वेच्छ संयोगांखाली संवृत असतात:
    प्रत्येक $i \in I$ साठी $U_i \in
    \Top{O}$ असेल, तर $\bigcup \Setabs{U_i}{i \in
      I} \in \Top{O}$.
  \end{enumerate}
\end{defn}

विवृत संचांचा संग्रह संदर्भावरून कळत असेल,
तर सांस्थितिक अवकाश दर्शवण्यासाठी फक्त $X$
लिहू शकतो. मात्र $X$ ला विवृत संचांचा संग्रह
दिल्यानंतरच त्याला सांस्थितिक अवकाश म्हणता येते.

\begin{defn}
  अंतःप्रज्ञावादी विधानीय तर्कशास्त्राचे
  \emph{सांस्थितिक प्रतिमान} म्हणजे
  $\mModel{X} = \tuple{X, \Top{O}, V}$ हे
  त्रिक, जिथे $\Top{O}$ ही~$X$ वरील
  संस्थिति आहे आणि $V$ हे प्रत्येक विधानीय
  चलाला~$\Top{O}$ मधील विवृत संच नेमणारे
  फलन आहे.

  सांस्थितिक प्रतिमान~$\mModel{X}$ दिले असता
  $\Prop{X}{!A}$ ची आगमनाने पुढीलप्रमाणे
  व्याख्या करता येते:
  \begin{enumerate}
  \item $\Prop{X}{\lfalse} = \emptyset$
  \item $\Prop{X}{p} = V(p)$
  \item $\Prop{X}{!A \land !B} = \Prop{X}{!A} \cap \Prop{X}{!B}$
  \item $\Prop{X}{!A \lor !B} = \Prop{X}{!A} \cup \Prop{X}{!B}$
  \item $\Prop{X}{!A \lif !B} = \Interior{(X \setminus \Prop{X}{!A}) \cup
    \Prop{X}{!B}}$
  \end{enumerate}
  येथे $\Interior{V}$ हे $V \subseteq X$
  या संचाला त्याच्या \emph{अंतर्भागाकडे} नेणारे
  फलन आहे; म्हणजे त्या संचात सामावलेल्या सर्व
  विवृत संचांचा संयोग. दुसऱ्या शब्दांत,
  \[
  \Interior{V} = \bigcup \Setabs{U}{U \subseteq V \text{ आणि } U \in \Top{O}}.
  \]
\end{defn}

कोणत्याही संचाचा अंतर्भाग हा विवृत संच असतो,
कारण तो विवृत संचांचा संयोग असतो. म्हणून
$\Prop{X}{!A}$ हा नेहमी विवृत संच असतो.

सांस्थितिक चिन्हार्थमीमांसा अत्यंत अमूर्त
असली, तरी तिच्यामागची प्रेरणा समजून घेण्याचे
मार्ग आहेत. $X$ चे !!{element}s, म्हणजे
“बिंदू”, ही अशी स्थळे आहेत की जिथे विधानांचे
मूल्यांकन करता येते, असे समजा. ज्या सर्व
बिंदूंवर $!A$ सत्य आहे त्यांचा संच म्हणजे~$!A$
ने व्यक्त केलेले विधान. बिंदूंचा कोणताही संच
संभाव्य विधान नसतो; फक्त $\Top{O}$ चे
!!{element}s तसे असतात. $!A \Entails !B$
तेव्हा आणि केवळ तेव्हाच $!A$ सत्य असलेल्या
प्रत्येक बिंदूवर $!B$ सत्य असते; म्हणजे
प्रत्येक~$X$ साठी $\Prop{X}{!A}
\subseteq \Prop{X}{!B}$. असंगत विधान~$\lfalse$
कधीच सत्य नसते; म्हणून
$\Prop{X}{\lfalse} = \emptyset$.

अंतःप्रज्ञावादी तर्कशास्त्रात वैध असणारे
नियम टिकावेत, म्हणजे सर्व सांस्थितिक अवकाशांसाठी
$!A \Proves !B$ तेव्हा आणि केवळ तेव्हाच
$\Prop{X}{!A} \subseteq \Prop{X}{!B}$
असावे, यासाठी $!B \land !C$,
$!B \lor !C$ आणि $!B \lif !C$ यांनी
व्यक्त केलेली विधाने $!B$ आणि~$!C$
यांनी व्यक्त केलेल्या विधानांशी कशी संबंधित
असावीत? कोणत्याही $!A$ साठी
$\lfalse \Proves !A$ अपेक्षित आहे; कारण
प्रत्येक $U$ साठी $\emptyset
\subseteq U$ असल्याने ते पूर्ण होते.
$!B \land !C \Proves !B$ असल्याने
$\Prop{X}{!B \land !C} \subseteq \Prop{X}{!B}$
ही अट हवी; तसेच
$\Prop{X}{!B \land !C} \subseteq \Prop{X}{!C}$.
$W \subseteq U$ आणि $W \subseteq V$
या अटी पूर्ण करणारा सर्वांत मोठा संच
$U \cap V$ आहे. उलट,
$!B \Proves !B \lor !C$ आणि
$!C \Proves !B \lor !C$, म्हणून
$\Prop{X}{!B} \subseteq \Prop{X}{!B \lor !C}$
आणि $\Prop{X}{!C} \subseteq \Prop{X}{!B \lor !C}$
या अटी हव्यात. $U \subseteq W$ आणि
$V \subseteq W$ असलेला सर्वांत लहान
संच~$W$ म्हणजे $U \cup V$.

$\lif$ साठीची व्याख्या अधिक सूक्ष्म आहे:
$!A \lif !B$ हे, $!A$ शी संयोग केल्यावर
$!B$ निष्पन्न होईल, असे सर्वांत दुर्बल
विधान व्यक्त करते. $!A \lif !B$ चा
$!A$ शी संयोग केल्यावर~$!B$ निष्पन्न होते,
हे $(!A \lif !B)
\land !A \Proves !B $ वरून स्पष्ट आहे.
म्हणून $\Prop{X}{!A \lif !B}$ हा
$\Prop{X}{!A \lif !B} \cap \Prop{X}{!A}
\subseteq \Prop{X}{!B}$ ही अट पूर्ण करणारा
सर्वांत मोठा विवृत संच असला पाहिजे;
यावरून आपली व्याख्या मिळते.

\end{document}
